\chapter{IoT Architecture}
\label{ch:iot}

\section{Communication Protocol}

The system uses MQTT v3.1.1 with Eclipse Mosquitto as the broker. Each ESP32 publishes a bundled JSON telemetry message to a single topic per device:

\begin{verbatim}
building/{device_id}/telemetry
\end{verbatim}

The payload contains timestamped measurements for load, PV, battery, and environment in a single JSON object, reducing MQTT overhead compared to per-measurement topics.

\section{Payload Schema}

\begin{lstlisting}
{
  "timestamp": "2026-09-15T14:30:00+01:00",
  "device_id": "esp32-metering-01",
  "load": {"voltage_v": 228.3, "current_a": 7.82,
           "power_w": 1785.0, "energy_kwh": 42.37,
           "pf": 0.98, "freq_hz": 50.02},
  "pv":   {"voltage_v": 36.2, "current_a": 8.1,
           "power_w": 293.2},
  "battery": {"voltage_v": 51.4, "current_a": 3.2,
              "soc_pct": 72.0},
  "env":  {"temperature_c": 34.1, "humidity_pct": 45.0,
           "irradiance_wm2": 820.0}
}
\end{lstlisting}

All payloads are validated on the backend using Pydantic schemas before database insertion, ensuring data integrity.

\section{Security}

The MQTT broker requires username/password authentication (\texttt{allow\_anonymous false}). No credentials are hard-coded; all secrets are stored in environment variables via a \texttt{.env} file excluded from version control.

\section{Quality of Service}

Telemetry data uses QoS~0 (at most once delivery), which is acceptable for periodic 15-second measurements where occasional packet loss does not affect the 15-minute aggregation. Alerts and commands use QoS~1 (at least once).
